\documentclass[10pt]{article}
\usepackage[T1]{fontenc}
\usepackage[utf8]{inputenc}
\usepackage[french]{babel}
\usepackage{amsmath,amssymb,mathrsfs,mathpazo}
\usepackage{geometry}
\geometry{a4paper,margin=2cm}
\usepackage{graphicx,array,tabularx,tabularray,colortbl,multirow,multicol,booktabs,enumitem,pifont,ulem,fancybox,diagbox,cancel}
\usepackage{pstricks,pstricks-add,pst-plot,pst-tree,pst-node,pst-text,pst-func,pst-math,pst-bspline,pst-3dplot,pst-eucl,pst-solides3d}
\let\pscalculate\undefined
\usepackage{pst-bezier}
\usepackage{tikz,pgfplots,tkz-tab}
\usetikzlibrary{arrows,arrows.meta,patterns,calc,positioning,trees,decorations.markings,shapes.geometric,shapes.misc,shapes.arrows,intersections,decorations.pathmorphing,decorations.pathreplacing,backgrounds,shadows,babel}
\pgfplotsset{compat=1.18}
\usepgfplotslibrary{fillbetween}
\usepackage[np]{numprint}
\usepackage{listings,makecell,pas-tableur,twoopt}
\let\barre\undefined
\usepackage[pstricks]{tablvar}
\definecolor{darkgreen}{RGB}{0,100,0}
\providecommand{\red}{\color{red}}
\providecommand{\blue}{\color{blue}}
\providecommand{\green}{\color{green}}
\providecommand{\black}{\color{black}}
\providecommand{\white}{\color{white}}
\providecommand{\gray}{\color{gray}}
\providecommand{\cyan}{\color{cyan}}
\providecommand{\magenta}{\color{magenta}}
\providecommand{\yellow}{\color{yellow}}
\providecommand{\violet}{\color{violet}}
\providecommand{\orange}{\color{orange}}
\providecommand{\euro}{\texteuro}
\providecommand{\EUR}{\texteuro}
\providecommand{\ieme}{\textsuperscript{e}}
\providecommand{\degres}{\ensuremath{^\circ}}
\providecommand{\degree}{\ensuremath{^\circ}}
\providecommand{\vv}[1]{\overrightarrow{#1}}
\providecommand{\Coord}[1]{#1}
\providecommand{\decosix}{}
\providecommand{\starredbullet}{\ensuremath{\bullet}}
\pagestyle{empty}
\setlength{\parindent}{0pt}
\usepackage{adjustbox}\DeclareUnicodeCharacter{2236}{:}\DeclareUnicodeCharacter{2190}{\ensuremath{\leftarrow}}
\providecommand{\R}{\mathbb{R}}
\providecommand{\N}{\mathbb{N}}
\providecommand{\Z}{\mathbb{Z}}
\providecommand{\Q}{\mathbb{Q}}
\providecommand{\D}{\mathbb{D}}
\providecommand{\C}{\mathbb{C}}
\providecommand{\e}{{\mathrm{e}}}
\providecommand{\vect}[1]{\overrightarrow{#1}}
\providecommand{\vectt}[1]{\overrightarrow{\mathrm{#1}}}
\providecommand{\vvt}[1]{\overrightarrow{\mathrm{#1}}}
\providecommand{\np}[2][]{{#2~#1}}
\providecommand{\vv}[1]{{\overrightarrow{#1}}}
\renewcommand{\d}{{\mathrm{d}}}
\providecommand{\degres}{{^\circ}}
\providecommand{\degre}{{^\circ}}
\providecommand{\blue}{\color{blue}}
\providecommand{\red}{\color{red}}
\providecommand{\enskip}{\quad}
\providecommand{\llbracket}{\lbrack\!\lbrack}
\providecommand{\rrbracket}{\rbrack\!\rbrack}
\providecommand{\euro}{{\text{€}}}
\providecommand{\no}{{\text{n°}}}
\providecommand{\interval}[2]{{[#1;#2]}}
\providecommand{\intervFO}[2]{{[#1;#2[}}
\providecommand{\dotfill}{\ldots}

\providecommand{\ds}{\displaystyle}
\providecommand{\textding}[1]{#1}
\providecommand{\Oij}{\(\mathrm{O};\vec{\imath},\vec{\jmath}\)}
\providecommand{\Oijk}{\(\mathrm{O};\vec{\imath},\vec{\jmath},\vec{k}\)}
\providecommand{\Ouv}{\(\mathrm{O};\vec{u},\vec{v}\)}
\providecommand{\intervFF}[2]{\ensuremath{[#1;#2]}}
\renewcommand{\euro}{\ifmmode\text{\texteuro}\else\texteuro\fi}
\DeclareUnicodeCharacter{00B0}{\ifmmode^{\circ}\else\textdegree\fi}
\DeclareUnicodeCharacter{20AC}{\ifmmode\text{\texteuro}\else\texteuro\fi}
% Énoncés : Éducation nationale. Transcription et corrigés : APMEP, Denis Vergès et contributeurs. Découpage et classement : CoopMaths, Sébastien Lozano, Éric Elter et contributeurs. Adaptation de la présentation : Maths Explorer.
% CC BY-SA : https://coopmaths.fr/www/about/
% https://coopmaths.fr/alea/static/bac/2025/tex/bac_2025_06_sujet2_polynesie_4.tex
\begin{document}
\noindent\textit{APMEP / CoopMaths. CC BY-SA. Présentation adaptée par Maths Explorer.}\par\medskip
\medskip

\emph{Pour chacune des affirmations suivantes, indiquer si elle est vraie ou fausse.\\
Chaque réponse doit être justifiée. Une réponse non justifiée ne rapporte aucun point.}

\medskip

\begin{enumerate}
\item Soient $E$ et $F$ les ensembles $E = \{1~;~2~;~3~;~4~;~5~;~6~;~7\}$ et 

$F = \{0~;~1~;~2~;~3~;~ 4~;~5~;~6~;~7~;~8~;~9\}$.

\smallskip

\textbf{Affirmation \no 1 :} Il y a davantage de 3-uplets d'éléments distincts de $E$ que de combinaisons à 4 éléments de $F$.
\end{enumerate}

\begin{minipage}{0.65\linewidth}
\begin{enumerate}[resume]
\item Dans le repère orthonormé ci-contre, on a représenté la fonction carré, notée $f$, ainsi que le carré ABCD de côté 3.

\textbf{Affirmation \no 2 :} La zone hachurée et le carré ABCD ont la même aire.
\end{enumerate}
\end{minipage}\hfill
\begin{minipage}{0.3\linewidth}
\psset{unit=0.5cm,arrowsize=2pt 3}
\begin{pspicture}(-4.5,-1)(4.5,10)
\psaxes[linewidth=1.25pt,labelFontSize=\scriptstyle]{->}(0,0)(-4.5,0)(4.5,10)
\psframe[linecolor=blue,linewidth=1.35pt](3,3)
\pscustom[fillstyle=vlines]{\psplot[plotpoints=2000,linewidth=1.25pt,linecolor=red]{0}{3.}{x dup mul}\psline(3,9)(3,0)(0,0)}
\uput[dl](0,0){\blue \small A}\uput[dr](3,0){\blue \small B}\uput[ur](3,3){\blue \small C}\uput[ul](0,3){\blue \small D}
\psplot[plotpoints=2000,linewidth=1.25pt,linecolor=red]{-3.1}{3.1}{x dup mul}
\end{pspicture}
\end{minipage}

\begin{enumerate}[start=3]
\item On considère l'intégrale $J$ ci-dessous : 

\[J = \displaystyle\int_1^2 x \ln (x)\,\text{d}x.\]

\textbf{Affirmation \no 3 :} Une intégration par parties permet d'obtenir: $J = \dfrac{7}{11}$.

\item Sur $\R$, on considère l'équation différentielle 

\[(E) :\quad  y' = 2y - \e^x.\]

\textbf{Affirmation \no 4 :} La fonction $f$ définie sur $\R$ par $f(x) = \e^x + \e^{2x}$ est solution de l'équation différentielle $(E)$.
\item  Soit $x$ donné dans [0~;~1[. On considère la suite $(u_n)$ définie pour tout entier naturel $n$ par:

\[u_n = (x - 1)\e^n + \cos (n).\]

\textbf{Affirmation \no 5 :} La suite $(u_n)$ diverge vers $-\infty.$
\end{enumerate}

\end{document}
